% ======================================================================
%  Appendix B --- Task-induced regularity and target matching
% ======================================================================
\section{Task-induced regularity and target matching}
\label{app:tm}

This appendix gives the full statements and proofs behind
Section~\ref{sec:contrib-tm}: the non-membership result, the transverse-kink
normal form and its Sobolev threshold, the higher-order matching defect, the RKHS
non-membership corollary, the quantitative approximation barrier with its
matching one-dimensional upper bound, and the target-matching regret guarantee.

\subsection{Loss of differentiability and non-membership}

\begin{proof}[Proof of Proposition~\ref{prop:tm-nonmember}]
Let $h=f-T$ and $g=\abs{h}$. Since $f\in C^1$ and $f(x_0)=T$, $h$ is
differentiable at $x_0$ with $h(x_0)=0$, so $h(x_0+tv)=t\nabla h(x_0)^\top v+o(t)$.
Choose $v$ with $\nabla h(x_0)^\top v\neq0$ (possible since $\nabla f(x_0)\neq0$).
Then the directional difference quotient of $g$ satisfies
$\lim_{t\to0^+}\frac{g(x_0+tv)-g(x_0)}{t}=\abs{\nabla h(x_0)^\top v}$ while
$\lim_{t\to0^-}=-\abs{\nabla h(x_0)^\top v}$; the two one-sided limits differ, so
$g$ is not differentiable at $x_0$. If $\Hk\subset C^1(\X)$ then every element of
$\Hk$ is differentiable everywhere; hence $g\notin\Hk$.
\end{proof}

\begin{remark}
The transversality $\nabla f(x_0)\neq0$ is necessary: if $\nabla f(x_0)=0$ the
level set may be tangent and $\abs{f-T}$ may stay smooth (e.g.\ $f(x)=x^2$,
$T=0$, giving $\abs{f-T}=x^2\in C^\infty$).
\end{remark}

\subsection{Besov threshold of a one-dimensional kink}

\begin{lemma}[Besov regularity of the model profile]
\label{app:lem-besov}
Let $\phi\in C_c^\infty(\R)$ with $\phi\equiv1$ near $0$. Then
$w(u)=\phi(u)\abs{u}$ satisfies $w\in H^s(\R)$ for all $s<3/2$ and
$w\notin H^{3/2}(\R)$; more precisely $w\in B^{3/2}_{2,\infty}(\R)$ but
$w\notin B^{3/2}_{2,2}(\R)=H^{3/2}(\R)$.
\end{lemma}

\begin{proof}
Since $|u|\in L^1_{\mathrm{loc}}(\R)$, its distributional derivatives are
well-defined. For any $\psi\in C_c^\infty(\R)$, integration by parts on
$(-\infty,0)$ and $(0,\infty)$ gives
\[
\langle D|u|,\psi\rangle
=-\int_{\R}|u|\psi'(u)\,du
=\int_{\R}\operatorname{sgn}(u)\psi(u)\,du.
\]
Thus $D|u|=\operatorname{sgn}(u)$ in the distributional sense. Similarly,
\[
\langle D\operatorname{sgn},\psi\rangle
=-\int_{\R}\operatorname{sgn}(u)\psi'(u)\,du
=2\psi(0),
\]
so $D^2|u|=2\delta_0$. Applying the distributional Leibniz rule to
$w=\phi|u|$, we obtain
\[
w'=\phi'|u|+\phi\operatorname{sgn}(u),\qquad
w''=2\phi''|u|+2\phi'\operatorname{sgn}(u)+2\phi(0)\delta_0
=2\delta_0+r,
\]
where $r\in C_c^\infty(\R)$, since $\phi\equiv1$ near the origin.

As $\widehat{\delta_0}=1$, the Fourier characterization of Sobolev spaces
yields
\[
\|\delta_0\|_{H^\sigma}^2
=\int_{\R}(1+\omega^2)^\sigma\,d\omega<\infty
\quad\Longleftrightarrow\quad \sigma<-\frac12.
\]
Since $r\in H^\sigma(\R)$ for every $\sigma\in\R$, it follows that
$w''\in H^\sigma(\R)$ if and only if $\sigma<-1/2$.
Moreover, $w\in L^2(\R)$, and the Fourier identity
$\widehat{w''}(\omega)=-\omega^2\widehat w(\omega)$ implies
\[
w\in H^s(\R)\quad\Longleftrightarrow\quad
w''\in H^{s-2}(\R).
\]
Consequently, $w\in H^s(\R)$ if and only if $s<3/2$.

For the sharp Besov regularity, consider the centered second difference
$\Delta_h^2w(u)=w(u+h)-2w(u)+w(u-h)$. Since $\phi\equiv1$ near
the origin, for sufficiently small $h$ and $|u|\le |h|$,
\[
\Delta_h^2w(u)=\Delta_h^2|u|
=2(|h|-|u|).
\]
Hence
\[
\|\Delta_h^2w\|_2^2
\ge 4\int_{-|h|}^{|h|}(|h|-|u|)^2\,du
=\frac83|h|^3.
\]
Conversely, the same local calculation, together with the smoothness
and compact support of $w$ away from the origin, gives
$\|\Delta_h^2w\|_2^2\le C|h|^3$ for sufficiently small $h$.
Thus
\[
\|\Delta_h^2w\|_2\asymp |h|^{3/2}.
\]
By the second-difference characterization of Besov spaces,
$w\in B^{3/2}_{2,\infty}(\R)$. On the other hand,
$B^{3/2}_{2,2}(\R)=H^{3/2}(\R)$, so the Sobolev result above
implies $w\notin B^{3/2}_{2,2}(\R)$.
\end{proof}


\begin{theorem}[Transverse kink imposes the $3/2$ threshold]
\label{app:thm-normalform-initial}
Let $\prm:\R^m\to\R$ be continuous of the form $\prm=\prm_+$ on $\{\Phi\ge0\}$ and
$\prm=\prm_-$ on $\{\Phi\le0\}$, with $\prm_\pm\in C^2$ near
$\Sigma=\{\Phi=0\}$ and distinct normal derivatives $\partial_n\prm_+\neq\partial_n\prm_-$
on $\Sigma$. Let $F\in C^2$ be transverse to $\Sigma$ at $x_0$. Then
$S=F^{-1}(\Sigma)$ is locally a $C^2$ hypersurface, and in suitable local $C^2$
coordinates $(u,z)\in\R\times\R^{d-1}$ straightening $S=\{u=0\}$,
\[
g(u,z)=h(u,z)+a(z)\abs{u},\qquad h\in C^2,\ a(z_0)\neq0 .
\]
Consequently $g\in H^s_{loc}$ for all $s<3/2$ and $g\notin H^{3/2}_{loc}$.
\end{theorem}

\begin{proof}
\emph{Straightening.} Since $\nabla\Phi\neq0$ on $\Sigma$ and $F$ is transverse,
$\psi:=\Phi\circ F\in C^2$ satisfies $\psi(x_0)=0$ and
$\nabla\psi(x_0)=DF(x_0)^\top\nabla\Phi(F(x_0))\neq0$. By the implicit function
theorem $S=\{\psi=0\}$ is a $C^2$ hypersurface, and $u:=\psi$ completed by $C^2$
coordinates $z$ along $S$ gives a local $C^2$ diffeomorphism whose sign
distinguishes the two branches. \emph{Normal form.} On each branch
$\prm_\pm\circ F\in C^2$; expanding to first order transversally,
$g=\prm_\pm\circ F=c(z)+b_\pm(z)u+O(u^2)$ ($\pm u>0$), with $c$ continuous
(because $\prm_+=\prm_-$ on $\Sigma$) and $b_+(z_0)\neq b_-(z_0)$ (distinct normal
derivatives). Setting $a(z)=\tfrac12(b_+(z)-b_-(z))$ and absorbing the regular
part (average of the two branches) into $h\in C^2$ yields
$g=h+a(z)\abs{u}$ with $a(z_0)\neq0$. \emph{Regularity.} Invariance of $H^s$ for
$s\le2$ under $C^2$ diffeomorphisms reduces the regularity of $g$ to the normal
form. The term $a(z)\abs{u}$ is treated by Lemma~\ref{app:lem-besov} in $u$; the
$C^2$ dependence in $z$ (with $a(z_0)\neq0$) and integration over the $d-1$
tangential variables preserve the critical exponent $3/2$ by Fubini for weighted
Sobolev norms, the singularity being purely normal. Hence $g\in H^s_{loc}$ for
$s<3/2$ and $g\notin H^{3/2}_{loc}$.
\end{proof}

\begin{theorem}[Transverse kink imposes the $3/2$ threshold]
\label{app:thm-normalform}
Let $\Phi\in C^2(\R^m)$ with $\nabla\Phi\neq0$ near
$y_0\in\Sigma:=\{\Phi=0\}$. Let $\prm:\R^m\to\R$ be continuous,
with $\prm=\prm_+$ on $\{\Phi\ge0\}$ and $\prm=\prm_-$ on
$\{\Phi\le0\}$, where $\prm_\pm\in C^2$ near $y_0$ and
$\partial_n\prm_+(y_0)\neq\partial_n\prm_-(y_0)$.
Let $F:\R^d\to\R^m$ be $C^2$, with $F(x_0)=y_0$, and
transverse to $\Sigma$ at $x_0$. Then $S=F^{-1}(\Sigma)$
is locally a $C^2$ hypersurface. In suitable local $C^2$
coordinates $(u,z)\in\R\times\R^{d-1}$ straightening
$S=\{u=0\}$, the function $g=\prm\circ F$ satisfies
\[
g(u,z)=c(z)+b(z)u+a(z)\abs{u}+O(u^2),
\qquad a(z_0)\neq0.
\]
Moreover, $g\in H^s_{\mathrm{loc}}$ for every $s<3/2$,
but $g\notin H^{3/2}(V)$ for any sufficiently small
neighborhood $V$ of $x_0$.
\end{theorem}

\begin{proof}
\emph{Straightening.} Set $\psi:=\Phi\circ F$. By transversality,
\[
\nabla\psi(x_0)
=DF(x_0)^\top\nabla\Phi(y_0)\neq0.
\]
The implicit function theorem implies that $S=\{\psi=0\}$
is locally a $C^2$ hypersurface. Taking $u:=\psi$ and
completing it with tangential coordinates $z$ yields a
local $C^2$ diffeomorphism that straightens $S=\{u=0\}$
and distinguishes the two branches by the sign of $u$.

\emph{Normal form.} Write $g_\pm=\prm_\pm\circ F$ in these
coordinates. Since $g_\pm\in C^2$ and the two branches
coincide on $S$, Taylor expansion in the transverse
direction gives
\[
g_\pm(u,z)=c(z)+b_\pm(z)u+O(u^2),
\qquad b_\pm(z)=\partial_u g_\pm(0,z).
\]
Setting $b=(b_++b_-)/2$ and $a=(b_+-b_-)/2$ yields
\[
g(u,z)=c(z)+b(z)u+a(z)\abs{u}+O(u^2).
\]
To verify that $a(z_0)\neq0$, let $q=\prm_+-\prm_-$.
Since $q=0$ on $\Sigma$, its gradient at $y_0$ is
normal to $\Sigma$. The distinct normal derivatives
therefore imply
$\nabla q(y_0)=\lambda\nabla\Phi(y_0)$ for some
$\lambda\neq0$. As $u=\Phi\circ F$, the chain rule gives
\[
b_+(z_0)-b_-(z_0)
=\partial_u(q\circ F)(0,z_0)=\lambda\neq0.
\]

\emph{Sobolev regularity.} The continuity of $g$ across
$\{u=0\}$ implies that its first distributional derivative
has no singular contribution. Its transverse derivative
has a jump $J(z):=b_+(z)-b_-(z)=2a(z)$, so
\[
D_u^2g=r(u,z)+J(z)\delta_0(u),
\]
where $r=\partial_u^2g_\pm$ on the respective half-spaces
and $r\in L^2_{\mathrm{loc}}$. Since the tangential
derivatives of the two branches agree on $\{u=0\}$,
their second tangential derivatives produce no Dirac
contribution. Consequently,
\[
\Delta g=R(u,z)+J(z)\delta_0(u),
\qquad R\in L^2_{\mathrm{loc}}.
\]

Let $\chi\in C_c^\infty$ be a cutoff supported in the
coordinate neighborhood and equal to one near $(0,z_0)$.
Writing $v=\chi g$, the distributional product rule gives
\[
\Delta v=R_\chi(u,z)+A(z)\delta_0(u),
\qquad R_\chi\in L^2,\quad
A(z)=\chi(0,z)J(z).
\]
Here $A$ is compactly supported, belongs to $L^2$, and
is not identically zero since $J(z_0)\neq0$.
Taking the Fourier transform in $(u,z)$, with dual
variables $(\xi,\zeta)$, yields
\[
\widehat{A(z)\delta_0(u)}(\xi,\zeta)
=\widehat A(\zeta).
\]
Hence, for any $\sigma\in\R$,
\[
\|A\delta_0\|_{H^\sigma}^2
=\int_{\R^{d-1}}|\widehat A(\zeta)|^2
\int_{\R}(1+\xi^2+|\zeta|^2)^\sigma
\,d\xi\,d\zeta.
\]
The inner integral is finite precisely when
$\sigma<-1/2$; in this range it is uniformly bounded
in $\zeta$. Since $A\in L^2$ is nonzero,
\[
A(z)\delta_0(u)\in H^\sigma
\quad\Longleftrightarrow\quad
\sigma<-\frac12.
\]
As $R_\chi\in L^2$, the same threshold holds for
$\Delta v$ when $\sigma<0$. Finally, since $v\in L^2$,
the Fourier characterization gives
\[
\|v\|_{H^s}^2\asymp
\|v\|_{L^2}^2+\|\Delta v\|_{H^{s-2}}^2.
\]
It follows that $v\in H^s$ for every $s<3/2$,
whereas $v\notin H^{3/2}$. The argument applies to
arbitrary local cutoffs, and the $C^2$ coordinate
change preserves the relevant Sobolev regularity.
Thus $g\in H^s_{\mathrm{loc}}$ for all $s<3/2$,
while $g\notin H^{3/2}(V)$ in every sufficiently
small neighborhood of $x_0$.
\end{proof}


\begin{corollary}[Target matching]
\label{app:cor-tm}
For $f\in C^2$, $f(x_0)=T$, $\nabla f(x_0)\neq0$, $g=\abs{f-T}$ satisfies
$g\in H^s_{loc}$ for $s<3/2$ and $g\notin H^{3/2}_{loc}$.
\end{corollary}
\begin{proof}
Take $m=1$, $\Phi(u)=u-T$, $\prm_\pm(u)=\pm(u-T)$: the normal derivatives are
$+1$ and $-1$, distinct, and $F=f$ is transverse where $\nabla f\neq0$. Apply
Theorem~\ref{app:thm-normalform} with $a\equiv1$.
\end{proof}

\begin{theorem}[Order-$q$ matching defect]
\label{app:thm-orderq}
Under the hypotheses of Theorem~\ref{app:thm-normalform} but with
$\prm_\pm\in C^{q+1}$ and normal derivatives coinciding up to order $q-1$ and
differing at order $q$ (and $F\in C^{q+1}$), the normal form becomes
$g(u,z)=h(u,z)+a(z)u_+^q$, $u_+=\max(u,0)$, $h\in C^{q+1}$, $a(z_0)\neq0$, and
$g\in H^s_{loc}$ for $s<q+\tfrac12$, $g\notin H^{q+1/2}_{loc}$.
\end{theorem}
\begin{proof}
Identical to Theorem~\ref{app:thm-normalform}, except that matching to order $q-1$
cancels the lower Taylor terms, leaving a jump at order $q$, i.e.\ a term
$a(z)u_+^q$. The profile $u_+^q$ has weak derivative of order $q$ equal to
$q!\,\mathbf{1}_{u>0}$, whose Fourier transform places it in
$H^{-1/2-\eta}\setminus H^{-1/2}$; integrating $q+1$ times gives the threshold
$q+1/2$.
\end{proof}

\begin{corollary}[RKHS non-membership]
\label{app:cor-nonmember}
Let $k$ be $\tau_k$-regular, $\Hk\simeq H^{\tau_k}$. Under the order-$q$ defect,
if $\tau_k\ge q+\tfrac12$ then $g\notin\Hk$; for a kink ($q=1$),
$\tau_k\ge3/2\Rightarrow g\notin\Hk$; for a Mat\'ern-$\nu$ kernel,
$\nu+\tfrac d2\ge q+\tfrac12\Rightarrow g\notin\Hk$.
\end{corollary}
\begin{proof}
If $g\in\Hk\simeq H^{\tau_k}_{loc}$ then $g\in H^{q+1/2}_{loc}$ once
$\tau_k\ge q+1/2$, contradicting Theorem~\ref{app:thm-orderq} (the large
inequality suffices, the critical case $\tau_k=q+1/2$ being itself excluded).
\end{proof}

\begin{remark}[Critical case and dimension]
$d$ is the domain dimension, not the number of objectives. In $d=1$ for a kink,
Mat\'ern-$1/2$ has $\tau_k=1<3/2$: a Lipschitz kink then \emph{belongs} to the
RKHS and the geometric misspecification disappears. The threshold is an
inequality; $\nu=1$ ($\tau_k=3/2$) sits exactly at the critical value and remains
misspecified. This does not make the task-in-model method correct in the presence
of noise: the distributional layer persists for every kernel.
\end{remark}

\subsection{Quantitative approximation barrier}

\begin{theorem}[Approximation barrier, restated]
\label{app:thm-barrier-initial}
Let $g=\abs{f-T}$ with $f\in C^2$, $f(x_0)=T$, $\nabla f(x_0)\neq0$, and $k$ a
$\tau_k$-regular kernel with $\nu:=\tau_k-\tfrac d2>1$, so that
$\Hk\hookrightarrow C^\nu$ with $\norm{h}_{C^\nu}\le C_S\Hnorm{h}$. Then there
exist $c>0$ and $B_0$ such that for all $B\ge B_0$,
$\epsinf(B)=\inf_{\Hnorm{h}\le B}\norm{h-g}_\infty\ge cB^{-1/(\nu-1)}$. In $d=1$,
mollification gives $\epsinf(B)\asymp B^{-1/(\nu-1)}$.
\end{theorem}

\begin{proof}
\emph{Lower bound.} Fix an even integer $r>\nu$ and a transverse direction $v$
(with $\partial_v(f-T)(x_0)\neq0$); let $\Delta_t^r$ be the order-$r$ centred
finite difference of step $t$ along $v$ at $x_0$. In normal-form coordinates
$g=h+a\abs{u}$ with $h\in C^2$, and, up to reduction, $a(z_0)\neq0$. A direct
computation of the even-order finite difference of $\abs{u}$ at $0$ gives
$\big|\Delta_t^r\abs{u}\big|_{u=0}=\kappa_r\,t$ with $\kappa_r>0$ (the even order
$r$ avoids the cancellation an odd centred difference would produce). Since
$h\in C^2$ and $r>2$, $|\Delta_t^r h|\lesssim t^2 t^{r-2}=t^r=o(t)$; hence for
small $t$,
\[
\big|\Delta_t^r g\big|_{x_0}\ge\tfrac12\abs{a(z_0)}\kappa_r\,t=:c_1 t .
\]
On the other hand, for $h\in\Hk$ with $\Hnorm{h}\le B$, the embedding
$\Hk\hookrightarrow C^\nu$ gives
$|\Delta_t^r h|\le C_r\norm{h}_{C^\nu}\le C_r C_S B\,t^\nu$ (a difference of order
$r>\nu$ of a $C^\nu$ function is $O(t^\nu)$). The operator $\Delta_t^r$ has
$\ell^1$-norm $2^r$, so
\[
2^r\norm{g-h}_\infty\ge\big|\Delta_t^r(g-h)\big|_{x_0}
\ge\big|\Delta_t^r g\big|_{x_0}-\big|\Delta_t^r h\big|_{x_0}
\ge c_1 t-C_r C_S B\,t^\nu .
\]
Optimizing the right-hand side over $t$ (maximum at $t^\star\asymp B^{-1/(\nu-1)}$)
gives value $\asymp B^{-1/(\nu-1)}$, whence $\norm{g-h}_\infty\gtrsim
B^{-1/(\nu-1)}$; taking the infimum over $h$ gives the bound locally near $x_0$,
which suffices since $\norm{g-h}_{\infty,\X}\ge\norm{g-h}_{\infty,U}$ for a
neighbourhood $U$ of $x_0$.

\emph{Upper bound ($d=1$).} Mollify $g_\delta=\varphi_\delta*g$ with a smooth
mollifier of scale $\delta$. Since $g$ is Lipschitz, $\norm{g-g_\delta}_\infty\le
C\delta$, while a standard estimate gives
$\norm{g_\delta}_{\Hk}\lesssim\delta^{-(\tau_k-1)}=\delta^{-(\nu-1/2)}$ in $d=1$.
Taking $B\asymp\delta^{-(\nu-1)}$ (after the fine optimization of
\S\ref{app:tm}, analogous to the RBF computation) yields
$\epsinf(B)\lesssim B^{-1/(\nu-1)}$, matching the lower bound.
\end{proof}


\begin{theorem}[Approximation barrier]
\label{app:thm-barrier}
Let $\mathcal X\subset\mathbb R^d$ contain a neighbourhood of $x_0$,
and let $g=|f-T|$, where $f\in C^2(\mathcal X)$,
$f(x_0)=T$, and $\nabla f(x_0)\neq0$.
Suppose that $\mathcal H_k$ embeds continuously into
$B^\nu_{\infty,\infty}(\mathcal X)$, with
$\nu:=\tau_k-d/2>1$. Then there exist $c>0$ and $B_0>0$
such that, for all $B\ge B_0$,
\[
\varepsilon_\infty(B)
:=\inf_{\|h\|_{\mathcal H_k}\le B}\|h-g\|_\infty
\ge cB^{-1/(\nu-1)}.
\]
In dimension $d=1$, if $\mathcal X$ is a compact interval,
$f$ is $C^2$ on a neighbourhood of $\mathcal X$ with finitely
many simple crossings of the level $T$, and
$\mathcal H_k$ is norm-equivalent to the restriction of
$H^{\tau_k}(\mathbb R)$, then the bound is sharp:
\[
\varepsilon_\infty(B)\asymp B^{-1/(\nu-1)}.
\]
\end{theorem}

\begin{proof}
\emph{Lower bound.}
Choose a unit vector $v$ such that
$a:=|\partial_v f(x_0)|>0$. Taylor's theorem and the
reverse triangle inequality give
\[
g(x_0+sv)=a|s|+O(s^2).
\]
Fix an even integer $r>\nu$ and let $\Delta_t^r$ denote
the centred finite difference of order $r$ and step $t$
along $v$. Since
$\Delta_t^r|s||_{s=0}=\gamma_r t$ with
$|\gamma_r|=\kappa_r>0$, we obtain
\[
\Delta_t^r g(x_0)=a\gamma_r t+O(t^2).
\]
Thus, for sufficiently small $t>0$,
$|\Delta_t^r g(x_0)|\ge c_1t$.

By the continuous Besov embedding, every
$h\in\mathcal H_k$ with $\|h\|_{\mathcal H_k}\le B$
satisfies
\[
|\Delta_t^r h(x_0)|\le C_rBt^\nu.
\]
Since the sum of the absolute values of the finite-difference
coefficients is $2^r$, it follows that
\[
2^r\|g-h\|_\infty
\ge |\Delta_t^r(g-h)(x_0)|
\ge c_1t-C_rBt^\nu.
\]
Choosing
$t=(c_1/(\nu C_rB))^{1/(\nu-1)}$,
which is admissible for sufficiently large $B$, yields
\[
\|g-h\|_\infty\ge cB^{-1/(\nu-1)}.
\]
Taking the infimum over the RKHS ball proves the lower bound.

\emph{Upper bound ($d=1$).}
Write $\tau=\tau_k=\nu+1/2>3/2$.
Under the additional assumptions, $g$ admits a compactly
supported extension $G$ to $\mathbb R$ whose second
distributional derivative is a finite signed measure.
Consequently,
\[
|\widehat G(\xi)|\le C|\xi|^{-2},
\qquad |\xi|\ge1.
\]
Let $G_\delta=\varphi_\delta*G$, where $\varphi$ is a
smooth mollifier. Since $G$ is Lipschitz,
\[
\|G-G_\delta\|_\infty\le C\delta.
\]
Moreover, using
$\widehat{G_\delta}(\xi)
=\widehat\varphi(\delta\xi)\widehat G(\xi)$,
we obtain
\[
\begin{aligned}
\|G_\delta\|_{H^\tau}^2
&\le C+
C\int_{|\xi|\ge1}
|\xi|^{2\tau-4}
|\widehat\varphi(\delta\xi)|^2\,d\xi\\
&\le C\delta^{-(2\tau-3)}.
\end{aligned}
\]
The last inequality follows by the change of variable
$u=\delta\xi$, since $\tau>3/2$ and
$\widehat\varphi$ decays rapidly.
By norm equivalence, the restriction
$h_\delta=G_\delta|_{\mathcal X}$ satisfies
\[
\|h_\delta\|_{\mathcal H_k}
\le C\delta^{-(\tau-3/2)}
=C\delta^{-(\nu-1)}.
\]
Taking $\delta=(C/B)^{1/(\nu-1)}$ gives
$\|h_\delta\|_{\mathcal H_k}\le B$ and
\[
\varepsilon_\infty(B)
\le\|g-h_\delta\|_\infty
\le C'B^{-1/(\nu-1)}.
\]
Together with the lower bound, this proves the claim.
\end{proof}


\begin{remark}[Locality and the RBF exclusion]
The lower bound is local, living exactly where the optimizer must be precise. The
RBF kernel is excluded from the polynomial statement because its RKHS is not a
finite-order Sobolev space; qualitatively an analytic kernel is even more strongly
penalized by a cusp, and the decay of $\epsinf(B)$ is then slower than any power.
A precise RBF lower bound is not claimed here.
\end{remark}

\subsection{Target-matching regret guarantee}

\begin{theorem}[Separation of concerns for target matching, restated]
\label{app:thm-tmsoc}
Under the confidence event of Theorem~\ref{app:thm-conf} with $f\in\Hk$,
$\Hnorm{f}\le B$, the rule $x_t\in\argmin_x a_t^{\mathrm{TM}}(x)$,
$a_t^{\mathrm{TM}}(x)=\abs{\mu_{t-1}(x)-T}-\sqrt{\beta_t}\sigma_{t-1}(x)$, gives,
with probability $\ge1-\delta$,
$R_T^{\mathrm{TM}}=\sum_{t=1}^T(\abs{f(x_t)-T}-\abs{f(x^\star)-T})
=O(\sqrt{T\beta_T\gain_T})$.
\end{theorem}

\begin{proof}
Write $g(x)=\abs{f(x)-T}$ and $r_t=g(x_t)-g(x^\star)$. By the triangle inequality
and the confidence bound at $x$,
\[
g(x)=\abs{f(x)-T}\le\abs{f(x)-\mu_{t-1}(x)}+\abs{\mu_{t-1}(x)-T}
\le\sqrt{\beta_t}\sigma_{t-1}(x)+\abs{\mu_{t-1}(x)-T}.
\tag{B.1}
\]
Symmetrically, $\abs{\mu_{t-1}(x)-T}\le\sqrt{\beta_t}\sigma_{t-1}(x)+g(x)$.
Applying (B.1) at $x_t$ and the definition of $x_t$
($a_t^{\mathrm{TM}}(x_t)\le a_t^{\mathrm{TM}}(x^\star)$, i.e.\
$\abs{\mu_{t-1}(x_t)-T}-\sqrt{\beta_t}\sigma_{t-1}(x_t)\le
\abs{\mu_{t-1}(x^\star)-T}-\sqrt{\beta_t}\sigma_{t-1}(x^\star)$), together with the
symmetric inequality at $x^\star$, gives, after rearrangement,
\[
g(x_t)\le\abs{f(x^\star)-T}+2\sqrt{\beta_t}\sigma_{t-1}(x_t)
=g(x^\star)+2\sqrt{\beta_t}\sigma_{t-1}(x_t),
\]
hence $r_t\le2\sqrt{\beta_t}\sigma_{t-1}(x_t)$, the same per-step bound as
classical GP-UCB. Squaring, summing, Cauchy--Schwarz and
Lemma~\ref{app:lem-sumvar} give $R_T^{\mathrm{TM}}=O(\sqrt{T\beta_T\gain_T})$.
\end{proof}

\noindent This is the $m=1$, $\prm(z)=-\abs{z-T}$ special case of the pushforward
theorem of Appendix~\ref{app:mo}.
